\documentclass[a4paper]{amsart}

\usepackage[usenames]{color}
\usepackage[colorlinks=true, citecolor=red]{hyperref} 
\usepackage{amssymb,latexsym}
\usepackage[alphabetic]{amsrefs}
\usepackage[mathscr]{eucal}
\usepackage{pdfsync}
\usepackage{graphicx}
\usepackage{epstopdf}
\DeclareGraphicsRule{.tif}{png}{.png}{`convert #1 `dirname #1`/`basename #1 .tif`.png}
\usepackage[all]{xy}
\usepackage{titletoc}

\CompileMatrices

\theoremstyle{plain}
\newtheorem{theorem}{Theorem}[section]
\newtheorem{lemma}[theorem]{Lemma}
\newtheorem{proposition}[theorem]{Proposition}
\newtheorem{corollary}[theorem]{Corollary}

\theoremstyle{definition}
\newtheorem{definition}[theorem]{Definition}%[section]

\theoremstyle{remark}
\newtheorem{example}[theorem]{Example}
\newtheorem{remark}[theorem]{Remark}
\newtheorem{discussion}[theorem]{Discussion}
\DeclareMathOperator{\Lip}{Lip}
\DeclareMathOperator{\dist}{dist}
\DeclareMathOperator{\expected}{\mathbb{E}}
\DeclareMathOperator{\prb}{Prob}
\DeclareMathOperator{\spt}{spt}
\DeclareMathOperator{\lap}{\Delta}

\raggedbottom

\addtolength{\textwidth}{2cm}
\addtolength{\oddsidemargin}{-1cm}
\addtolength{\evensidemargin}{-1cm} 
\addtolength{\topmargin}{-1cm} 
\addtolength{\textheight}{2cm}
% \renewcommand{\l@subsection}{\@contentsline{2}{5em}{7em}}

%\makeatletter
%\makeatother

 
%%% BEGIN DOCUMENT
\begin{document}

\title{Markov Processes}  
\author{John E.\ Hutchinson}
\date{\today}

\maketitle


%\renewcommand*\l@subsection{\@tocline{4}{50em}{7em}{10em}{10em}}


\tableofcontents

\newpage \part{Types of Markov Processes}

I follow mainly \cite{Lam}*{Chapter  6}.

%%%%%%%%%%%%%%%%%%%%%%%%%%%%%%%%%%%%%%
%%%%%%%%%%%%%%%%%%%%%%%%%%%%%%%%%%%%%%%%
%%%%%%%%%%%%%%%%%%%%%%%%%%%%%%%%%%%%%%%%%%
\section{\textbf{Continuous Time and Discrete Space}}

Assume states  $ \{1,\dots, N\}$. 

\subsection{Markov transition function $ P_{ij}(t) $ }
The process $ X(t) $ should represent the state at time $ t \geq 0 $.
$ P(t) $ is \emph{interpreted} as the $ N \times N $ matrix such that 
\begin{equation} \label{} 
P_{ij} (t) = P(X(t) = j \mid X(t) = i )\quad t \geq 0  .
\end{equation}


 
Motivated by this,  a \emph{Markov transition function } is defined to be a set $ (P(t))_{t\geq 0} $  of   $ N \times N $ matrices   s.t.
\begin{align}
& P_{ij}(t) \geq 0\ \forall i,j, \quad \sum\nolimits_j  P_{ij}(t) = 1 \ \forall i, \quad \text{i.e.\ $ P(t) $ is a \emph{stochastic} matrix},  \label{P1}\\
&P(0) = I_n, \label{P2}\\
&P(t+s) = P(t) P(s) \quad \forall s,t \geq 0, \label{P3}\\
& P(t) \text{ is right continuous at $ t =0 $}. \label{P4}
\end{align} 
 
%%%%%%%%%%%%%%%%%%%%%
\subsection{Exponential}  The \emph{exponential} of a square matrix $ A $ is 
\begin{equation} \label{}
e^A = \exp{A} := \sum_{n\geq 0 }\frac{A^n}{n!}. 
\end{equation}
The series converges since $ A_{ij} \leq \|A\|^n $.  Moreover, it is straightforward to check 
\begin{align}
& e^{A+B} = e^A e^B \quad \text{if } AB = BA,\\
& \frac{d}{dt}e^{tA} = Ae^{tA} = e^{tA} A .
\end{align}


%%%%%%%%%%%%%%%%%%%%%
\subsection{Generator} \label{secgen}
\begin{theorem}
Suppose $ (P(t))_{t\geq 0} $ is a Markov transition function.

Then  $ P(t) $ is differentiable, and in particular continuous, for all $ t $.  Let 
\begin{equation} \label{defG}
G = \left. \frac{d}{dt}\right|_{t=0+} P(t)= \lim_{t\to 0+}\frac{P(t) - I}{t}.
\end{equation}
It follows 
\begin{align} 
&P(t) = e^{tG} \quad t\geq 0,   \label{dG1}\\
& G_{ij}\geq 0 \text{ for } i\neq j , \quad \sum_jG_{ij}= 0 \ \forall i. \label{dG2}
\end{align} 
Moreover, 
\begin{equation} \label{diffP}
P'(t) = P(t)\, G = G\, P(t). 
\end{equation}
 
 Conversely, if $ G $ satisfies \eqref{dG2} and  $ P(t) $ is defined by \eqref{dG1} then $ P(t) $ is a Markov transition  function. 
 
 The correspondence between $ P $ and $G $ is one-one.
\end{theorem}

\begin{proof}\mbox{}

\textbf{1}: \emph{Continuity} of $ P(t) $ for all $ t $ follows from \eqref{P3} and~\eqref{P4}.

\medskip
\textbf{2}: To see $ P(t) $ is \emph{everywhere differentiable} note
\begin{gather*}
 \big(P(h)-I\big)\big( I + P(h) + \cdots + P((n-1)h)\big) = P(nh) - I,\\
 \therefore \frac{P(h)-I}{h}\, h \big( I + P(h) + \cdots + P((n-1)h)\big)  =  P(nh) - I.
 \end{gather*}
 
 Fix $ t>0 $, let $ h \to 0 $ and $ n\to \infty $ s.t.\ $ nh \to t $.
Then%
\footnote{From the theorem, the integral on the right side is
\[   \int_0^t e^{uG}\, du = G^{-1} (e^{tG}-I) = G^{-1} (P(t) - I).
\]
}
\[
\lim_{h\to 0^+}  \frac{P(h)-I}{h}  = (P(t) - I) \left( \int_0^t P(u)\, du\right)^{-1},
\]
using continuity of $ P $ and the Riemann sum approximation to the integral.

\medskip \textbf{3}:  Define $ G $ as in \eqref{defG}.  Then \eqref{dG2} is clear. Moreover $ P'(t) $ exists for all $ t $ and satisfies \eqref{diffP} since  $ P'(0) $ exists and by~\eqref{P3}. 

Since $ P(t) $ satisfies \eqref{diffP}, and by standard ODE theory the unique solution which equals $ I $ at $ =0 $ is  $ e^{tG} $, it follows \eqref{dG1} is true.  

\medskip
\textbf{4}: Finally suppose $ G $ is as in  \eqref{dG2}.  Then we claim $ P(t) := e^{tG} $ is a Markov transition matrix.  Then \eqref{P2}, \eqref{P3} and \eqref{P4} are clear.  

For \eqref{P1}  we first note the rows of $ P(t) $ sum to $ 1 $ since
\[
e^{tG} \boldsymbol{1} = I \boldsymbol{1} = \boldsymbol{1},
\]
because the rows of $ G$  sum to $ 0 $.

It remains to prove $ P(t) \geq O$, the matrix of zeros. 

First assume $ G_{ij}>0 $ for $ i\neq j $.  Then
for small $ t> 0 $ 
\[
P(t) =  I + t G + o(t) >0.
\]
Hence for \emph{arbitrary}  $ t>0 $, $ P(t) = (P(t/k))^k >O$, by taking  $ k $ sufficiently large.

If we only have $ G \geq 0 $, let $E_{\epsilon}   $ be the matrix with all entries $ \epsilon > 0 $. Then 
\[  P = \lim_{\epsilon \to 0} e^{G + E_{\epsilon} } \geq O, \] 
since $ e^{G + E_{\epsilon} } > O$ by the previous argument.
\end{proof} 



%%%%%%%%%%%%%%%%%%%%%%%%%%%%%%%%%%%%%%
\subsection{Random Clock} \label{secrc} The natural probabilistic way to generate   a continuous time Markov chain $ (P(t)_{t\geq 0} $ is via a discrete time Markov chain $ Q $ governed by a random clock following a Poisson R.V.\  with parameter $ \lambda $.

\medskip More precisely, transitions occur at times given by a Poisson R.V.\ with mean $ \lambda t $ in time intervals of length $ t $.
Each transition is governed by a transition matrix $ Q $.

Thus the probability of $ k $ transitions in $ [0,t] $ is $ e^{- \lambda t } ( \lambda t )^k / k! $, and the probability of passing from state $  i $ at time $ 0 $ to state $ j $ by time $ t $ is
\begin{equation} \label{}
P_{ij}(t) = e^{-\lambda t }  \sum_{k\geq 0} \frac{(\lambda t)^k}{k!} Q_{ij}^k.
\end{equation}
Equivalently, 
\begin{equation} \label{}
P(t) = e^{\lambda t (Q-I)}.
\end{equation}
In particular, \emph{the generator of $ (P(t))_{t\geq 0 } $ is $ \lambda (Q-I) $}.

\medskip Conversely, if $ P(t) $ is a given transition function with generator $ G $  then we can describe the process in the above manner by taking  
\begin{equation} \label{}
Q = I + \lambda ^{-1} G, \quad \lambda \geq \max_i |G_{ii} |.
\end{equation}
The reason for the restriction on $ \lambda $ is so that  $ 0 \leq Q_{ii} \leq 1 $.  Recall $ -1 \leq G_{ii} \leq 0 $
from~\eqref{dG2}.

%%%%%%%%%%%%%%%%%%%%%%%%%%%%%%%%%%%%%%
%%%%%%%%%%%%%%%%%%%%%%%%%%%%%%%%%%%%
%%%%%%%%%%%%%%%%%%%%%%%%%%%%%%%%%%%%%%%%%%
\section{\textbf{Discrete Time and General State Space}}



\subsection{Markov kernel $ P(x,E) $}

Consider a measurable state-space $ ( S, \mathcal{S} ) $.

$ P(x,E) $  is \emph{interpreted} as \emph{the probability of moving from $ x $ into $ E $ in one transition.}

Motivated by this, a \emph{Markov kernel} or \emph{stochastic kernel} is defined to be a function $ P(x,E) $ where $ x\in S $ and $  E\in \mathcal{S} $, such that  $ P(x,E) $ is a probability measure on $   ( S, \mathcal{S} ) $ for fixed $ x $,  and is measurable in $ x $ for fixed $ E $.

%%%%%%%%%%%%%%%%%%%%%%f
\subsection{Properties}
Define
\begin{equation} \label{dtg1}
\begin{aligned}
P^0(x,E) &= \delta _x, \\
P^1(x,E) &= P(x,E), \\
P^{n+1} (x,E) &= \int P(x,dy) P^n(y,E).
\end{aligned}
\end{equation}
Then $ P^n $ gives the $n  $-step probabilities.

It follows readily that
\begin{equation} \label{dtg2}
P^{n+m} (x,E) = \int P^n(x,dy)\, P^m(y,E).
\end{equation} 
 

%%%%%%%%%%%%%%%%%%%%%%%%%%%
\subsection{Redistributions}  \label{secred} Instead of  probabilities or probability distributions $ P(x,\cdot) $, it is often helpful to think of a mass distribution $ \nu_x $ and its redistributions.

Define
\begin{equation} \label{rdm} 
\begin{aligned}
&\nu_x(E)  = P(x,E), \text{ i.e.\  $ \nu_x $ is the  redistribution of $ \delta_x $ after one step},\\
&\nu_{n,x}(E)   = P^n(x,E), \text{   $ \nu_{n,x}  $ is the  redistribution of $ \delta_x $ after $ n $  steps},\\
&\nu_\mu(E)   = \int P(x,E)\, d\mu(x)= \int \nu_x(E)\, d\mu(x) , \text{  $ \nu_\mu  $ is the  redistribution of $\mu $ after one  step}. 
\end{aligned} 
\end{equation} 
Note 
\[
P(x,dy) =   d\nu_x(y), \quad P^n(x,dy) =   d\nu_{n,x}(y).
\]
Then  \eqref{dtg2}  becomes
\[  \nu_{m+n,x}(E) = \int d\nu_{n,x}(y)\, \nu_{m,y}(E). \]
That is, the redistribution of $ \delta_x $ after $ n+m $ steps equals the redistribution after $ n $ steps \emph{of} the redistribution of $ \delta_x $  after $ m $ steps.




%%%%%%%%%%%%%%%%%%%%%%%%%%%%%%%%%%%%%%
%%%%%%%%%%%%%%%%%%%%%%%%%%%%%%%%%%%%
%%%%%%%%%%%%%%%%%%%%%%%%%%%%%%%%%%%%%%%%%%
\section{\textbf{Continuous Time and State Space}}\label{secctss}

\subsection{Markov transition function $ P_t(x,E) $}  Consider a measurable state-space $ ( S, \mathcal{S} ) $.

$ P_t(x,E) $ is \emph{interpreted as the probability that a system in state $ x $ at time $ s $ will move to some state in $ E $ at time $ s+ t $}.

\medskip Motivated by this, a \emph{Markov transition function} is defined to be a function $ P_t(x,E) $  where $ t\geq 0 $, $ x\in S $ and $ E \in \mathcal{S} $.  The requirements are that $ P_t(x,E) $ be a probability measure on $   ( S, \mathcal{S} ) $ for fixed $ t $ and $ x  $,  and measurable in $ x $ for fixed $ t $ and $ E $. Moreover, we require
\begin{equation} \label{mtf}
\begin{aligned}
&P_0(x,E) = \mathcal{X} _E(x)  \quad (\text{i.e.\ } P_0(x,\cdot) = \delta_x),\\
& P_{t+s}(x,E) = \int_S P_t(x,dy)\,  P_s(y,E) \quad \text{if } s,t\geq 0 .
\end{aligned} 
\end{equation}
The last equation is the \emph{Chapman Kolmogorov} equation and reflects the Markov principle of lack of ``memory''.

\medskip
In terms of \emph{mass distributions} as in Section~\ref{secred},  set 
\[  \nu_{t,x} (E) = P_t(x,E). \]
Then the  Chapman Kolmogorov  equation becomes
\[  \nu_{t+s,x}(E) = \int d\nu_{t,x}  (y)\, \nu_{s,y}(E),  \]
and says that the redistribution of $ \delta_x $ at time $ t+s $  equals the redistribution at   time $ t $ of the redistribution of $ \delta_x $  at time $ s $.




%%%%%%%%%%%%%%%%%%%%%%%%%%%%%%%%%
\subsection{Brownian motion}
Suppose a probability measure $ \mu $ is \emph{infinitely divisible}.%
\footnote{That is, for each integer $ n>0 $ there is a probability measure $ \mu_n $ for which 
\[ \mu = \mu_n * \cdots * \mu_n \text{\ $ n $ times}.\]
 Recall $ \mu*\nu = \nu*\mu $ and 
\[ \mu*\nu(E) = \iint \mathcal{X}_{E}(x+y)\, d\mu(x)\, d\nu(y) 
   =   \int \mu(E-y) \, d\nu(y).
   \]
   Recall also that if $ \dist  X  = \mu $, $ \dist  Y  =\nu $, and $ X $ and $ Y $ are independent, then $ \dist (X+Y) = \mu * \nu $.
   }
It follows that for each $  t>0 $ there is a probability measure $ \mu_t $ such that $ \mu_1=\mu $ and  $ \mu_{s+t} = \mu_s * \mu_t $.  

It follows that for $ S = \mathbb{R} $ and $ \mathcal{S} = \mathcal{B} $ (Borel sets on $ \mathbb{R} $),
\[
P_t(x,E) :=   \mu_t (E-x) 
\]
defines a translation invariant measure. That is $ P_t(x, E) = P_t(x+y, E+y)$.  The proof is essentially notational.%
\footnote{One has for $ t,s > 0 $, 
 \begin{align*}
 P_{t+s} &= \mu_{t+s}(E-x) = \mu_s * \mu_t (E-x) \\
 & = \int \mu_s(E-x-y)\, d\mu_t(y)  = \int d\mu_t(y-x)\,\mu_s(E-y)
 \intertext{($ y \mapsto y-x $, and where $ \mu_t(y-x) $ is translation of $ d\mu_t$ to right by $ x $)}
 &=\int P_t(x,dy)\, P_s(y,E).
 \end{align*}
For this to be true at $ t=0 $ it is clearly necessary and sufficient that $ P_0(x,\cdot) = \delta_x $.
}






It follows that the process has \emph{stationary independent increments}.%
\footnote{\cite{Lam}*{p125} says ``independent increments''.  But this is true for any Markov process.}

\medskip  The most important example is $ \mu = N(0,1) $, the normal distribution with mean 0 and variance~1.%
\footnote{Recall that if $ \dist X = N(0,s) $ and $ \dist X = N(0,t) $, and $ X $ and $ Y $ are independent, then 
$ \dist (X+Y) = N(0,s+t) $.}
Then $ \mu_t = N(0,t) $ (variance $ t $), i.e.
\begin{equation} \label{} 
P_t(x,E) = \frac{1}{\sqrt{2\pi t}} \int_{E-x} e^{-y^2/2t} dy.
\end{equation} 
This is the transition function for \emph{Brownian motion} [BM].

The \emph{continuity condition} 
\begin{equation} \label{} 
P_t(x, [x-\epsilon, x+\epsilon ] ) = 1 - o(t) \quad \epsilon > 0 
\end{equation} 
is clearly satisfied.  
This implies the process moves by continuous motion.  Such a markov process is called a \emph{diffusion}.%
\footnote{If one defines a markov process by using a stochastic kernel $ Q(x,E) $ and a random clock analogous to the manner  in Section~\ref{secrc} then  one instead obtains a \emph{jump} process.}

 

%%%%%%%%%%%%%%%%%%%%%%%%%%%%%%%%%%%%%%
%%%%%%%%%%%%%%%%%%%%%%%%%%%%%%%%%%%%%%%%
%%%%%%%%%%%%%%%%%%%%%%%%%%%%%%%%%%%%%%%%%%
\newpage \part{Semigroups}

I follow mainly \cite{Lam}*{Chapter  7}.


\section{\textbf{Construction of Semigroups from Markov Processes}}   
 
\subsection{Notation}
$ P_t(x,E) $ is a Markov transition function on $ (S, \mathcal{S} ) $ as in Section~\ref{secctss}. In particular note~\eqref{mtf}. 


 
Define the Banach spaces
\begin{align} 
F(S) & = \{ f:S \to \mathbb{R} \mid   \text{ is bdd and measurable} \} \text{ with the sup norm $ \|f\| $},\label {bsl} \\
C(S)  & = \{ f:S \to \mathbb{R} \mid   \text{ is bdd and continuous} \} \text{ with the sup norm $ \|f\| $},\label {csl} \\
M(S) &= \{ \mu \mid \text{$ \mu $ is a finite signed measure on $ S $}\} \text{ with the variation norm $ \|\mu\| $}, 
\label{bsm}
\end{align} 
where   $ \|\mu\| = \mu^+(S) + \mu^-(S) $. 

 
%%%%%%%%%%%%%%%%%%%
\subsection{Contraction semigroups on $ F(S) $ and $ M(S) $} (This is a `warm up''. The important case is Section~\ref{seccs}.)

Define 
\begin{align}
&T_t:F(S)\to F(S), \quad  T_tf (x) = \int f(y) P_t(x,dy), \\
& U_t:M(S) \to M(S), \quad  U_t\mu(E) = \int d\mu(x)\, P_t(x,E) =  \int d\mu(x)\, \nu_{t,x}(E).\label{sgu}
\end{align} 
Thus if $ X_t $ is the process corresponding to $ P(x,E) $ then we have the interpretation
\[ T_t f(x) = \mathbb{E} ^x f(X_t ), \]
Analogous to  the third equation in \eqref{rdm},  we  interpret  $ U_t \mu $ as the redistribution $ \nu_{t,\mu} $ of $ \mu $ at time~$ t $.  

Since $  U_t\mu $ is a measure it operates on $ f\in F (S) $, and using dual pairing notation we have
\begin{align*}
\langle U_t\mu  , f \rangle &:= U_t\mu(f) \\
   & =    \int  d\mu(x) \left( \int P_t(x,dy)\, f(y) \right)  \text{ since from \eqref{sgu} we have the case $ f = \mathcal{X}_E$}\\
   &= \int d\mu(x)\, T_t f(x) = \langle \mu, T_t f\rangle .
\end{align*}                           
That is,
\begin{equation} \label{}
\langle U_t\mu  , f \rangle =  \langle \mu, T_t f\rangle .
\end{equation} 
This can be interpreted for $ \mu $ a probability measure as follows: The mean of $ f $ w.r.t.\ the redistributed $ \mu $ equals the mean obtained by choosing $ x $ via $ \mu $, starting the process from $ x $ and evaluating $ f $ at $ X_t $, i.e.\ evaluating $ \mathbb{E}  \big( \mathbb{E} ^x f(X_t) \big)$.

\medskip It is easy to check%
\footnote{In particular,
\begin{align*} 
T_{t+s}f(x) & = \int P_{t+s}(x, dy ) \, f(y)\text{ by definition}\\
& = \iint P_t(x,du)\, P_s(u,dy) f(y) \text{ by the Chapman Kolmogorov eqn \eqref{mtf}}\\
&= \int P_t(x,du) T_sf(u) = (T_t\circ T_s f)(x) \text{ by definition}.
\end{align*}
To show $ U_t $ is contractive split $ \mu $ into $ \mu^+ $ and $ \mu^-$.  Inequality may occur as some of the redistribution of $ \mu^+ $ and  $ \mu^- $ may cancel.
}
that $ T_t  $ and $ U_t $ define contraction semigroups, i.e.
\begin{align}
 &\|T_t f \| \leq \|f\|, \quad T_{t+s} = T_t \circ T_s, \\
 &\|U_t \mu \| \leq \|\mu\|, \quad U_{t+s} = U_t \circ U_s.
 \end{align} 
 
\bigskip
In order to have a useful theory, we will need the \emph{continuity condition} which, in the case of $ F(S) $, is 
\[
\lim_{t\to 0^+} \| T_t f - f \| = 0  \quad \forall f\in F(S).
\]
This is equivalent to%
\footnote{Straightforward, see \cite{Lam}*{pp 154,5}.}  
\[\lim_{t\to 0^+}  P_t(x, \{x\} ) = 1 \quad \forall x\in S .\]

However, this is not true for B.M.   


%%%%%%%%%%%%%%%%%%%%%%%%
\subsection{Contraction semigroups  on $ C(S) $} \label{seccs}

 The Markov transition function $ P_t(x,E) $ is \emph{Feller} if
  \begin{equation} \label{}
T_t : C(S) \to C(S) ,
\end{equation}
where $ C(S) $ is the space of bounded continuous functions with the sup norm.




This is the same as  requiring weak convergence of measures:%
\footnote{Since 
\[
(T_tf)(x) = \int P_t(x,dy)\, f(y)\, dy = P_t(x,\cdot)(f).
\]}
\begin{equation} \label{} 
P_t(x,\cdot) \rightharpoonup  P_t(x_o,\cdot) \text{ as } x\to x_0, \quad \forall t\geq 0.
\end{equation} 

Note there is no continuity requirement here on $ t $, which we now consider.

\bigskip \emph{Henceforth we assume $ S $ is  compact}.  

Processes on $ \mathbb{R} $ and $ \mathbb{R}^n $ will be studied by considering the one point compactifications $   \mathbb{R}^* $ and~$ \mathbb{R}^{n*} $.

\bigskip The transition function $ P_t(x,E) $ is \emph{uniformly stochastically continuous} if 
\begin{equation} \label{usc} 
\lim_{t\to 0^+} P_t(x, B_\epsilon(x)) = 1 \quad \forall \epsilon >0, \text{ uniformly for } x \in S.
\end{equation} 
It is straightforward to check%
\footnote{See \cite{Lam}*{PP 157,8}. The point is that since $ S $ is compact, $ f $ is uniformly continuous.} 
that this implies 
\begin{equation} \label{} 
\lim_{t\to 0^+}  \| T_t f - f \| = 0, \quad \forall f \in C(S) .
\end{equation} 
  
\bigskip The transition function $ P_t(x,E) $ is  \emph{normal}%
\footnote{Non standard notation from \cite{Lam}.}
if it is Feller and uniformly stochastically continuous. 

\bigskip \emph{Remark}:  For a Feller transition function on a compact state space, stochastic continuity (i.e.\ \eqref{usc} without the uniform requirement in $ x $) implies uniform stochastic continuity and hence normality.%
\footnote{See \cite{Lam}*{pp 159--161}.} 
 
\bigskip BM on $ \mathbb{R}^n$ is normal (acting on the  compactified space).
 
%%%%%%%%%%%%%%%%%%%%%%%%%%%%%%%%%
%%%%%%%%%%%%%%%%%%%%%%%%%%%%%%%%%
%%%%%%%%%%%%%%%%%%%%%%%%%%%%%%%%%
\section{\textbf{General Contraction Semigroups}}

\subsection{Definitions}
Assume $ (V,\|\cdot\|) $ is a real Banach space.  

Examples to keep in mind for the following are \eqref{bsl} and \eqref{bsm}.

\medskip A family $ (T_t)_{t\geq 0} $ of (bounded) linear operators $ T_t:V \to V $ is a \emph{contractive semigroup} if 
\begin{equation} \label{}
\begin{aligned} 
&\|T_t x \| \leq \| x \| \ \forall x \in V,\\
& T_0 = I,\ T_{t+s} = T_t \circ T_s \ \forall t,s \geq 0,\\
&\lim_{t\to 0^+}\| T_t x - x \| =0\ \forall x \in V.
\end{aligned} 
\end{equation}



The last condition is equivalent to continuity at $ 0 $  in the \emph{strong} topology. This is weaker than continuity in the \emph{uniform} topology, which normally does not hold.  
But in this context strong continuity   is equivalent to continuity in the \emph{weak} topology.

It also follows that $ T_t $ is \emph{uniformly strongly continuous}%
\footnote{That is
\[
\forall x\ \forall \epsilon>0 \ \exists \delta = \delta(x,\epsilon) \text{ s.t. }
t \in [0,\infty),\ |s|< \delta \Longrightarrow  \|T_{t+s}x -T_t x \| < \delta .
 \]
}
for $ t\in [0,\infty) $. The proof is straightforward from the semigroup property.%
\footnote{See \cite{Lam}*{pp 137,138.}}

%%%%%%%%%%%%%%%%%%%%%%%%%%%%%%%
\subsection{Infinitesimal generator} The infinitesimal generator $ A $ of $ (T_t)_{t \geq 0} $ is  defined by
\begin{equation} \label{}
Ax = \lim_{t\to 0^+} \frac{Tx-x}{t}.
\end{equation} 
The \emph{domain} $ \mathcal{D} _A $ of $ A $ is the set of $ x \in V$ for which the limit exists.

See Section~\ref{secgen}, where $ A $ is the matrix $ G $.

\begin{proposition}
For any contractive semigroup $ (T_t)_{t\geq 0} $,  $ \mathcal{D}_A $ is dense in $ V $.
\end{proposition}

\begin{proof} Suppose $ x \in V $.

The idea is to first show that if   $ f(s) $ is $  C^1 $ with compact support in $ (0,\infty) $, then 
\[
y(x) := \int f(s) T_sx \, ds \in \mathcal{D} _A 
\]
Then  take a sequence of such $ f_n(s) \geq 0$ with  
unit integral converging to $ \delta_0 $.

\medskip The first claim follows from
\begin{align*} 
Ay  &= \lim_{t\to 0^+} \frac{T_t y - y}{t}  \\
 &=  \lim_{t\to 0^+}\int f(s) \frac{T_{t+s}x -T_sx}{t} ds  \\
 &= \lim_{t\to 0^+} \int \frac{f(s-t)-f(s)}{t} T_s x\, ds \\
& =\int f'(s) T_s x \, ds \quad \text{since   $ f $ is $ C^1 $}.
\end{align*} 
Hence $ y \in \mathcal{D}_A$.  


\medskip Now take $ f_n $ with supports in $ (0,1/n) $ converging to $ \{0\} $,  $ \int f_n =1 $, $ f_n \geq 0 $.  Then
\[
\|y_n(x) - x \| = \left\| \int f_n(s) (T_s x-x ) \, ds \right\|  
 \leq \int f_n(s) \|T_s x - x \| \, ds \to 0.
\]
Here we used the fact $ T_s x \to x $ as $ s\to 0 $.
\end{proof} 

\bigskip
\emph{We now make the crucial observation that  $ T_t x $ satisfies a first order linear ODE}.

\medskip In the following the limit giving the derivative is in the norm topology.  This is always the case unless o.w.\ stated.

\begin{proposition}
If $ x \in \mathcal{D}_A $ then   $ T_t x  \in \mathcal{D}_A $ for all $ t>0 $.  Moreover,
\begin{equation} \label{tde}
\frac{d}{dt} T_t x = A (T_t x) = T_t (A x). 
\end{equation}
\end{proposition}

The proof is standard using the semigroup property.%
\footnote{See \cite{Lam}*{pp 139,140}.}

%%%%%%%%%%%%%%%%%%%%%%%%%%%%%%
\subsection{Bounded infinitesimal generator}
 If $ A $ is bounded (this is \emph{not} usually the case) then $ x \in \mathcal{D}_A $    for \emph{every} $ x $ , and by standard methods
\[  e^{tA }x :=\left( \sum_{n\geq 0} \frac{t^n A^n }{n!}\right) x  \]
is well defined.  Moreover, this is the \emph{unique} solution%
\footnote{Suppose $ y(t) $ is any solution.  Then  we compute
\[ \frac{d}{dt} \big(e^{-tA}y(t)\big) = 0, \quad    e^{-0A}y(0) = x.\]
It follows that $ e^{-tA}y(t)= x $ for all $ t\geq 0 $.  

The last ``follows'' is because $   u(t) \in V$ continuous on $ [0,a) $ with zero derivative on $ (0,a) $ implies $ u(t) = u(0) $ on $ [0,a] $.  To see this consider  arbitrary elements $ v$  of $ V^* $, differentiate  $ \langle u(t)-u(0), v \rangle $ and apply the standard real valued result    to deduce that  $  \langle u(t)- u(0), v \rangle = 0 $ for all $ v \in V^* $, and hence $ u(t) = u(0) $.}
of 
\[    \frac{dy(t)}{dt}    = A y(t), \quad y(0) =x.\]
It then follows from  \eqref{tde} that
\[  T_t x = e^{tA} x \quad  \forall t\geq 0 \text{ and } \forall x, \qquad T_t = e^{tA} \quad  \forall t \geq 0 .\]


%%%%%%%%%%%%%%%%%%%%%%%%%%%%%%
\subsection{Laplace transform analogy}

If the infinitesimal generator $ A $ is unbounded then $ e^{tA} x $ need not converge, even for $ x \in \mathcal{D}_A $.

So we try to solve the linear ODE \eqref{tde} by modelling the approach on the Laplace transform method.

\bigskip
Recall that for bounded real valued functions $ x:\mathbb{R}^+ \to \mathbb{R}  $ the Laplace transform is
\[
(\mathcal{L}x)(\lambda)= X(\lambda) := \int_0^\infty e^{-\lambda t} x(t)\, dt, \quad \lambda \geq 0. \]
Then differentiation and integration become multiplication and division by $ \lambda $.

Integration by parts gives
\[
(\mathcal{L} x')(\lambda) = \lambda X(\lambda) - x(0).
\]
The inverse transform we need is 
\[
X(\lambda ) = \frac{1}{\lambda - a } \ (\lambda > a) \Longleftrightarrow  x(t) = e^{-at}\ (t \geq 0 ).
\] 

\bigskip So if we want to solve the following ODE (the analogue of~\eqref{tde}) 
\[
\frac{dx  (t)}{dt} = ax(t), \quad x(0) = x_0,
\]
then translating this into an algebraic equation for the Laplace transform we get
\begin{align*}
\lambda X(\lambda) - x_0 &= a X(\lambda) \\
\therefore\ X(\lambda) &= (\lambda-a)^{-1} x_0 \quad \lambda >0 \\
\therefore\  x(t) &= e^{-at} x_0 \quad t\geq 0.
\end{align*} 

\bigskip  The above will motivate the approach to solving \eqref{tde} by using resolvents.
%%%%%%%%%%%%%%%%%%%%%%%%%%%%%%%%%%%%%%
%%%%%%%%%%%%%%%%%%%%%%%%%%%%%%%%%%%%
%%%%%%%%%%%%%%%%%%%%%%%%%%%%%%%%%%%%%%%%%%
\begin{thebibliography}{9} 


  
\bib{Lam}{book}{
   author={Lamperti, John},
   title={Stochastic processes},
   note={A survey of the mathematical theory;
   Applied Mathematical Sciences, Vol. 23},
   publisher={Springer-Verlag},
%   place={New York},
   date={1977},
   pages={xvi+266},
%   review={\MR{0461600 (57 \#1585)}},
}




\end{thebibliography}





\end{document}


 